% numodel-plot-example-quadrants.tex — axis placement across all
% nine quadrant combinations.
%
% Demonstrates how numodel-plot positions the axis lines and labels
% based on which quadrants the data range covers:
%   - 4 single-quadrant cases (I, II, III, IV)
%   - 4 half-plane cases (I+II, II+III, III+IV, IV+I)
%   - 1 four-quadrant case
% In every "axisMoved" case (the latter five) the axis label is
% placed in the extension of the axis arrow, 2em past the tip,
% rather than at the tick-label margin.

\documentclass{article}
\usepackage[a4paper, margin=1.5cm]{geometry}
\usepackage{siunitx}
\usepackage{caption}
\usepackage{subcaption}
\usepackage{numodel-plot}

% Compact plots so 9 fit on a page.
\numodelplotsetup{xcmmax=4, ycmmax=4}

% Shorthand: one plot of the given function over the supplied range.
\newcommand{\qplot}[5]{%
  \def\xmin{#1}\def\xmax{#2}%
  \def\ymin{#3}\def\ymax{#4}%
  \def\xlabelqty{t}\def\xlabelunit{\s}%
  \def\ylabelqty{v}\def\ylabelunit{\m\per\s}%
  \drawplot{\addplot[domain=\xmin:\xmax,thick,samples=40]{#5};}%
}

\begin{document}

\begin{figure}[h]\centering
% Row 1: single-quadrant cases
\begin{subfigure}{0.32\textwidth}\centering
  \qplot{0}{10}{0}{5}{0.5*x}
  \caption*{Quadrant I}
\end{subfigure}\hfill
\begin{subfigure}{0.32\textwidth}\centering
  \qplot{-10}{0}{0}{5}{-0.5*x}
  \caption*{Quadrant II}
\end{subfigure}\hfill
\begin{subfigure}{0.32\textwidth}\centering
  \qplot{-10}{0}{-5}{0}{0.5*x}
  \caption*{Quadrant III}
\end{subfigure}

\vspace{1em}
% Row 2: Q IV + the two horizontal half-planes
\begin{subfigure}{0.32\textwidth}\centering
  \qplot{0}{10}{-5}{0}{-0.5*x}
  \caption*{Quadrant IV}
\end{subfigure}\hfill
\begin{subfigure}{0.32\textwidth}\centering
  \qplot{-5}{10}{0}{5}{0.05*x*x}
  \caption*{I + II}
\end{subfigure}\hfill
\begin{subfigure}{0.32\textwidth}\centering
  \qplot{-10}{0}{-3}{3}{0.6*(x+5)}
  \caption*{II + III}
\end{subfigure}

\vspace{1em}
% Row 3: the two remaining half-planes + 4Q
\begin{subfigure}{0.32\textwidth}\centering
  \qplot{-5}{10}{-5}{0}{-0.05*x*x}
  \caption*{III + IV}
\end{subfigure}\hfill
\begin{subfigure}{0.32\textwidth}\centering
  \qplot{0}{10}{-3}{3}{0.6*x-3}
  \caption*{IV + I}
\end{subfigure}\hfill
\begin{subfigure}{0.32\textwidth}\centering
  \qplot{-5}{10}{-3}{5}{0.5*x}
  \caption*{All four}
\end{subfigure}
\end{figure}

\end{document}
